% ============================================================
% APPENDIX -- KL-CAP BINDING, PAIR-LEVEL POOLING AND THE TWO SCALES (written 2026-11-06, review revision)
% Self-contained \subsection (app:kl-cap; tab:kl-cap-subsets, tab:lofo, tab:scales), \input from sections/08_appendix.tex
% after 08f, so it continues the Additional Results appendix. Answers reviewer concerns A1, A2, B1 (data part), C1 and C5
% of notes/REVIEW_v5_reviewer_agent_20261106.md. Every number comes from results/analysis/review_sel_summary.md and the
% review_sel_*.csv files it lists (written by scripts/analysis/review_selection.py, 2026-11-06T23:44Z); the src comment
% after each paragraph or table names the file.
% ============================================================
\subsection{KL-Cap Binding, Pair-Level Pooling and the Two Scales}
\label{app:kl-cap}

\paragraph{KL-cap binding.} The selection rule caps first-token KL divergence
at 0.1 for every model. The cap binds in a run when some trial above 0.1 had
fewer dev keyword refusals than the selected trial; a run stayed below the cap
when at most one of its trials exceeded 0.1. In the main study the cap bound
in 15 of the 24 direction-based runs (7 of 12 older and 8 of 12 newer models)
and in all 8 ARA runs, every one of which reached 0 dev refusals above the cap
(maximum KL 10.9 to 33.2). Eight direction-based runs stayed below the cap,
four on each generation: Qwen3.8-27B MPOA, Olmo-3.1-32B MPOA and SOM and
Mistral-Small-2409 MPOA never sampled a trial above 0.1, and OLMo-2-32B DiM,
MPOA and SOM and Qwen3.8-27B SOM sampled one. On the small pair the cap bound
in 4 of the 16 direction-based runs of both seeds (all four J-space runs
stayed below it) and in all 4 ARA runs. The selected dev KL ranged from 0.018
to 0.100 for the direction-based runs (median 0.080) and from 0.049 to 0.099
for ARA; 7 of 24 direction-based and 1 of 8 ARA selections lay within 0.01 of
the cap.
% src: results/analysis/review_sel_kl_summary.csv (main seed 1 direction: 15 of 24, old 7/12, new 8/12, left the cap at most once
% src: 4 + 4; optimisation 8 of 8, max_kl 10.90-33.25; small seed 1 and seed-2 direction: 4 of 16 binding, 8 of 16 at most once;
% src: small ARA 4 of 4; median_sel_kl 0.0796; n_sel_kl_within_0p01_of_cap 7 and 1) = review_sel_summary.md part 1;
% src: review_sel_kl_runs.csv (never_left_cap, left_cap_at_most_once; sel_kl range 0.018-0.100 direction, 0.049-0.099 ARA).

\paragraph{Position in the search range.} The selected trial's parameters were
read from the Optuna journals and placed in Heretic's default ranges
(Table~\ref{tab:search}); a maximum weight counts as at the ceiling when it
lies in the top 5\% of its range. The selected trial was at the ceiling in 5
of the 24 main-study direction-based runs and in none of the 8 runs that
stayed below the cap. In those 8 runs the selected maximum weights lay at a
median of 0.81 (attention; range 0.56--0.93) and 0.89 (MLP; 0.83--0.94) of
their range, against 0.58 and 0.84 in the other direction-based runs, so the
selected edits lean to the strong half of the range but not to its bound. The
search did sample the strongest maximum weights in each of the 8 runs, and
those trials had no fewer dev refusals than the selected trial in 7 of 8 (in
OLMo-2-32B SOM the better trial at the ceiling had KL 0.110). Counting every
strength-type parameter (maximum weight, minimum weight, minimum-weight
distance, SOM decay), the selected trial had at least one at the strong end of
its range in 14 of the 24 direction-based runs and in 4 of the 8 that stayed
below the cap (a minimum weight or minimum-weight distance in all four). The selected
trial was one of the 30 random start-up trials of the TPE sampler in 7 of the
24 direction-based runs (4 of the 8 that stayed below the cap) and in 1 of 8
ARA runs. Because seed~1 and the same search space give identical start-up
draws, two pairs of selected parameter vectors coincide across models
(OLMo-2-32B SOM and Qwen3.8-27B SOM; Olmo-3.1-32B MPOA and Qwen3.8-27B MPOA).
% src: results/analysis/review_sel_kl_params.csv, review_sel_kl_runs.csv (sel_max_weight_at_ceiling 5 of 24, 0 of 8;
% src: sel_attn_max_weight_pos / sel_mlp_max_weight_pos medians; n_trials_max_weight_at_ceiling, min_refusals_max_weight_at_ceiling;
% src: sel_is_startup_trial 7 of 24, 4 of 8, 1 of 8; sel_params_identical_to) = review_sel_summary.md part 1.

\paragraph{H1 and H2 by binding.} At the level of pair $\times$ method units,
the cap bound in both cells for Qwen DiM and ARA, all four Gemma methods, OLMo
ARA and Mistral Small DiM and ARA; in neither cell for Qwen MPOA and OLMo
MPOA; and in one cell for Qwen SOM, OLMo DiM and SOM, and Mistral Small MPOA
and SOM. Table~\ref{tab:kl-cap-subsets} refits the pooled GEE of
\S\ref{sec:res-h1} on these subsets. The raw direction-based effect was
positive in every subset; the baseline-adjusted one was negative where the
cap bound in both cells and positive where it did not, and the pooled +0.03 of
all units lies between the two. The subsets also differ in family, so they
cannot separate the cap from the family. H2 can be estimated only on units in
which the cap bound, because it bound in every ARA cell.
% src: results/analysis/review_sel_kl_units.csv (unit_category per pair x method); review_sel_kl_h1h2.csv; review_sel_summary.md
% src: part 1 (family composition of the subsets).

\begin{table}[t]
\centering\footnotesize
\setlength{\tabcolsep}{2.5pt}
\caption{Pooled H1 and H2 terms (GEE, log-odds, prompt clusters, family fixed effects, 95\% intervals) refitted on
subsets of the 16 main-study pair $\times$ method units, by whether the KL cap bound in each unit's two cells. Units:
direction-based / ARA. Dir.: direction-based class; adj.: baseline-adjusted (minus the originals' difference); n/a:
no ARA unit in the subset.}
\label{tab:kl-cap-subsets}
\begin{tabular}{@{}lcllll@{}}
\toprule
Subset & Units & Dir., raw & Dir., adj. & ARA, raw & H2 \\
\midrule
All units & 12/4 & +0.62 [+0.57, +0.68] & +0.03 [$-$0.05, +0.11] & +0.12 [+0.00, +0.24] & $-$0.50 [$-$0.62, $-$0.38] \\
Cap bound in both cells & 5/4 & +0.32 [+0.24, +0.41] & $-$0.24 [$-$0.33, $-$0.14] & +0.11 [+0.01, +0.22] & $-$0.21 [$-$0.33, $-$0.09] \\
Not bound in $\ge$1 cell & 7/0 & +1.03 [+0.95, +1.11] & +0.22 [+0.12, +0.32] & n/a & n/a \\
Bound in neither cell & 2/0 & +3.23 [+2.90, +3.56] & +2.20 [+1.95, +2.45] & n/a & n/a \\
$\ge$2 trials above the cap in both & 6/4 & +0.33 [+0.24, +0.41] & $-$0.28 [$-$0.37, $-$0.18] & +0.12 [+0.01, +0.23] & $-$0.21 [$-$0.33, $-$0.08] \\
$\le$1 trial above the cap in a cell & 6/0 & +1.04 [+0.96, +1.12] & +0.27 [+0.17, +0.37] & n/a & n/a \\
\bottomrule
\end{tabular}
\end{table}
% src: results/analysis/review_sel_kl_h1h2.csv (S0--S5; gee_planned and gee_with_originals terms) = review_sel_summary.md part 1 table.

\paragraph{Pair-level pooling.} With the pair as the unit, between-pair
heterogeneity dominated every pooled estimand ($I^2$ 99--100\%; between-pair
SD of the adjusted direction-based effect 22.9~pp and of the H2 contrast
26.1~pp). The 95\% prediction intervals for a new pair, which use a $t$
distribution on 2 degrees of freedom, were $-$5.90 to +6.02 log-odds for the
adjusted direction-based effect and $-$11.32 to +9.21 for the H2 contrast; on
the percentage-point scale they extend beyond $\pm$100~pp, so with four pairs
they state only that an effect of either sign is plausible for a new pair.
Mapping the percentage-point $\sigma$ to the logit scale by the delta method
at each cell's rate, instead of estimating $\sigma$ on the logit scale, gave
seed SEs up to 14.9 log-odds near 0\% refusal (Gemma-4-31B ARA) and pooled
intervals of $-$1.61 to +1.82 (adjusted direction-based) and $-$3.59 to +2.81
(H2). Table~\ref{tab:lofo} gives the leave-one-family-out GEE terms.
% src: results/analysis/review_sel_pair_meta.csv (I2, tau; pi_lo / pi_hi; delta-method variant rows "prompt + seed (pp sigma,
% src: delta method; unstable near 0%)": direction adjusted HKSJ [-1.610, 1.822], H2 log-odds [-3.594, 2.815]) = review_sel_summary.md
% src: part 2 (tau 22.9 pp / 26.1 pp; seed SE up to 14.9 log-odds for Gemma-4 ARA).

\begin{table}[t]
\centering\footnotesize
\setlength{\tabcolsep}{3pt}
\caption{Leave-one-family-out pooled GEE terms (log-odds; main study, WildGuard refusal, 972 harmful prompts).}
\label{tab:lofo}
\begin{tabular}{@{}lrrrrr@{}}
\toprule
Term & All four & w/o Qwen & w/o Gemma & w/o OLMo & w/o Mistral Small \\
\midrule
Generation effect, direction class & +0.62 & +0.69 & +0.81 & +0.24 & +0.74 \\
Generation effect, ARA & +0.12 & $-$1.37 & +1.02 & +0.24 & +0.15 \\
Generation $\times$ class (H2) & $-$0.50 & $-$2.06 & +0.20 & $-$0.00 & $-$0.59 \\
Baseline-adjusted, direction class & +0.03 & +0.22 & +0.13 & $-$0.51 & +0.33 \\
Baseline-adjusted, ARA & $-$0.46 & $-$1.68 & +0.31 & $-$0.51 & $-$0.25 \\
Originals' generation difference & +0.58 & +0.45 & +0.66 & +0.74 & +0.40 \\
\bottomrule
\end{tabular}
\end{table}
% src: results/analysis/review_sel_pair_lofo.csv (final.analysis_g; reproduces final_pooled_robustness.csv, max |difference|
% src: 3.7e-06) = review_sel_summary.md part 2(d).

\paragraph{The two scales.} Table~\ref{tab:scales} gives every pair $\times$
method generation effect, the originals' difference and the baseline-adjusted
effect on both scales. Raw effects of a single method have the same sign on
both scales, because they compare two rates. Adjusted effects changed sign
between scales in 3 of 21 cells and in significance in 5. On the percentage-point
scale all four Mistral Small adjusted effects were large and negative
($-$22.2 to $-$14.4~pp), whereas in log-odds they ranged from $-$1.37 to
$-$0.10.
% src: results/analysis/review_sel_scales.csv = review_sel_summary.md part 3 (0 of 21 raw effects differ; adjusted: sign 3 of 21,
% src: significance 5; Mistral adjusted -22.2 to -14.4 pp, log-odds -1.37 to -0.10).

\begin{table*}[t]
\centering\scriptsize
\setlength{\tabcolsep}{3pt}
\caption{Generation effect, originals' difference and baseline-adjusted generation effect per pair and method on the
percentage-point and the log-odds scale (WildGuard refusal, 972 harmful prompts; 95\% intervals, prompt sampling only).
Sign: the adjusted effect has opposite signs on the two scales. Sig.: its significance differs between scales.}
\label{tab:scales}
\begin{tabular}{@{}lllrllrlcc@{}}
\toprule
Pair & Method & Effect, pp & Orig., pp & Adjusted, pp & Effect, log-odds & Orig., log-odds & Adjusted, log-odds & Sign & Sig. \\
\midrule
Qwen & DiM & +26.9 [+23.8, +29.8] & +13.1 & +13.7 [+10.1, +17.1] & +1.16 [+1.02, +1.30] & +0.99 & +0.17 [$-$0.03, +0.37] &  & $\times$ \\
Qwen & MPOA & +10.1 [+7.8, +12.5] & +13.1 & $-$3.1 [$-$6.3, +0.2] & +1.39 [+1.04, +1.73] & +0.99 & +0.39 [+0.01, +0.77] & $\times$ & $\times$ \\
Qwen & SOM & $-$4.7 [$-$8.0, $-$1.4] & +13.1 & $-$17.8 [$-$21.5, $-$14.2] & $-$0.20 [$-$0.33, $-$0.06] & +0.99 & $-$1.19 [$-$1.40, $-$0.97] &  &  \\
Qwen & ARA & +34.4 [+31.2, +37.5] & +13.1 & +21.3 [+17.1, +25.4] & +2.35 [+2.08, +2.62] & +0.99 & +1.36 [+1.05, +1.68] &  &  \\
\midrule
Gemma & DiM & +16.4 [+13.5, +19.3] & +1.0 & +15.3 [+13.0, +17.8] & +1.18 [+0.96, +1.40] & +0.10 & +1.07 [+0.90, +1.25] &  &  \\
Gemma & MPOA & $-$4.7 [$-$8.6, $-$0.9] & +1.0 & $-$5.8 [$-$9.7, $-$1.7] & $-$0.19 [$-$0.35, $-$0.04] & +0.10 & $-$0.30 [$-$0.55, $-$0.04] &  &  \\
Gemma & SOM & +2.6 [$-$1.1, +6.2] & +1.0 & +1.5 [$-$1.9, +4.9] & +0.12 [$-$0.05, +0.28] & +0.10 & +0.01 [$-$0.22, +0.24] &  &  \\
Gemma & ARA & $-$26.0 [$-$28.9, $-$23.3] & +1.0 & $-$27.1 [$-$30.3, $-$23.6] & $-$4.07 [$-$4.87, $-$3.27] & +0.10 & $-$4.17 [$-$5.00, $-$3.34] &  &  \\
\midrule
OLMo & DiM & +34.8 [+31.5, +37.9] & +0.8 & +34.0 [+30.5, +37.6] & +1.45 [+1.31, +1.60] & +0.06 & +1.39 [+1.22, +1.56] &  &  \\
OLMo & MPOA & +61.4 [+58.2, +64.4] & +0.8 & +60.6 [+57.1, +64.1] & +3.91 [+3.56, +4.25] & +0.06 & +3.85 [+3.49, +4.21] &  &  \\
OLMo & SOM & +14.4 [+11.5, +17.3] & +0.8 & +13.6 [+10.4, +16.7] & +0.64 [+0.51, +0.78] & +0.06 & +0.58 [+0.42, +0.75] &  &  \\
OLMo & ARA & $-$2.5 [$-$4.6, $-$0.3] & +0.8 & $-$3.3 [$-$5.8, $-$0.9] & $-$0.23 [$-$0.43, $-$0.03] & +0.06 & $-$0.29 [$-$0.52, $-$0.07] &  &  \\
\midrule
Mistral Small & DiM & $-$4.3 [$-$7.4, $-$1.2] & +17.9 & $-$22.2 [$-$26.1, $-$18.4] & $-$0.20 [$-$0.35, $-$0.06] & +0.80 & $-$1.00 [$-$1.18, $-$0.82] &  &  \\
Mistral Small & MPOA & $-$2.9 [$-$4.7, $-$1.2] & +17.9 & $-$20.8 [$-$24.0, $-$17.7] & $-$0.57 [$-$0.92, $-$0.22] & +0.80 & $-$1.37 [$-$1.74, $-$1.00] &  &  \\
Mistral Small & SOM & +3.5 [+1.7, +5.3] & +17.9 & $-$14.4 [$-$17.7, $-$11.1] & +0.70 [+0.34, +1.06] & +0.80 & $-$0.10 [$-$0.47, +0.28] &  & $\times$ \\
Mistral Small & ARA & $-$0.7 [$-$2.1, +0.6] & +17.9 & $-$18.6 [$-$21.6, $-$15.7] & $-$0.17 [$-$0.45, +0.12] & +0.80 & $-$0.96 [$-$1.28, $-$0.65] &  &  \\
\midrule
Small pair & DiM & +34.1 [+30.5, +37.5] & +16.4 & +17.7 [+13.5, +21.8] & +1.54 [+1.36, +1.71] & +1.09 & +0.45 [+0.22, +0.68] &  &  \\
Small pair & MPOA & +4.2 [+1.6, +6.9] & +16.4 & $-$12.1 [$-$15.7, $-$8.3] & +0.32 [+0.12, +0.52] & +1.09 & $-$0.77 [$-$1.03, $-$0.51] &  &  \\
Small pair & SOM & +9.7 [+6.9, +12.4] & +16.4 & $-$6.7 [$-$10.4, $-$3.0] & +0.70 [+0.50, +0.90] & +1.09 & $-$0.39 [$-$0.65, $-$0.13] &  &  \\
Small pair & J-space & +21.2 [+18.1, +24.3] & +16.4 & +4.8 [+1.3, +8.3] & +0.93 [+0.79, +1.07] & +1.09 & $-$0.16 [$-$0.35, +0.04] & $\times$ & $\times$ \\
Small pair & ARA & +14.2 [+11.6, +16.9] & +16.4 & $-$2.2 [$-$6.0, +1.6] & +1.55 [+1.23, +1.86] & +1.09 & +0.46 [+0.10, +0.82] & $\times$ & $\times$ \\
\bottomrule
\end{tabular}
\end{table*}
% src: results/analysis/review_sel_scales.csv = review_sel_summary.md part 3 table (pp effects and adjusted pp with the paper's
% src: Newcombe / bootstrap intervals; log-odds by the delta method with the paired item covariance; flags "sign differs (adj.)" and
% src: "significance differs (adj.)").
